\section{The framework}\label{sec:framework}

Let $K$ be a field, let $V$ be a geometrically integral $K$-variety, and let
$\mathcal{H}$ be a countable family of height-bounded subsets of $V(K)$. Write
$\mathcal{H}(B)$ for the members of height at most $B$.

\begin{definition}\label{def:admissible}
A datum $x\in\mathcal{H}(B)$ is \emph{admissible} if its localization $x_v$
lies in a prescribed open subset $U_v\subseteq V(K_v)$ for every
$v\in\Omega$.
\end{definition}

We use the following model hypothesis to distinguish the structural input from
the analytic one.

\begin{assumption}\label{ass:equidistribution}
There are constants $\delta>0$ and $\eta>0$ such that, uniformly over
admissible residue conditions,
\begin{equation}\label{eq:counting}
  \#\mathcal{H}(B)=\delta B+O\!\left(B^{1-\eta}\right).
\end{equation}
\end{assumption}

\begin{remark}
In an application, $\delta$ is normally a product of local densities. The
substantive task is to prove that the product is nonzero and that the error
in~\eqref{eq:counting} is uniform.
\end{remark}
